% Generated draft: requires mathematical and source review.
\documentclass[11pt]{article}
\usepackage[margin=25mm]{geometry}
\usepackage{fontspec}
\setmainfont{Latin Modern Roman}
\usepackage{amsmath,amssymb,mathtools,unicode-math}
\setmathfont{Latin Modern Math}
\usepackage{enumitem,needspace,xurl,hyperref}
\setlength{\parindent}{0pt}
\setlength{\parskip}{5pt}
\setlength{\emergencystretch}{3em}
\newcommand{\N}{\mathbb{N}}
\newcommand{\Z}{\mathbb{Z}}
\newcommand{\Q}{\mathbb{Q}}
\newcommand{\R}{\mathbb{R}}
\newcommand{\C}{\mathbb{C}}
\newcommand{\F}{\mathbb{F}}
\newcommand{\A}{\mathbb{A}}
\newcommand{\PP}{\mathbb{P}}
\newcommand{\E}{\mathbb{E}}
\newcommand{\eps}{\varepsilon}
\newcommand{\dd}{\,\mathrm{d}}
\begin{document}
\section*{UnsolvedMath ID: 10600053}
\textbf{Source problem code:} catalogue entry
\subsection{Definitions and mathematical statement}
Let \(\mathcal G\) be the finite decorated graphs admitted as graph-link diagrams, and let \(\sim\) be equivalence generated by the graph-link Reidemeister moves. The supplied record does not specify the decorations, allowed graph class, or the exact moves, so these must be recovered from the intended graph-link theory.

\emph{Decision problem and conjecture.} Is there an algorithm which, on every input pair \((G,H)\in\mathcal G^2\), halts and correctly decides whether \(G\sim H\)? The source conjectures that there is no such algorithm. In a fixed effective encoding, this asks whether the predicate
\[
 \operatorname{Eq}(G,H)=
 \begin{cases}1,&G\sim H,\\0,&G\not\sim H\end{cases}
\]
is computable on all valid diagram pairs. The question concerns a total decision procedure for equivalence, not merely a method that sometimes recognizes equivalent pairs, a practical efficiency bound, or recognition of one distinguished trivial class.

The source offers heuristics, not a reduction establishing undecidability: arbitrary graphs and graph Reidemeister moves might encode Turing-machine computations, formal-language rules, or relations as in universal semigroups and finitely presented groups. It gives a rough analogy between decidable settings, naming low-dimensional topology and hyperbolic groups, and undecidable settings, naming higher-dimensional topology, arbitrary finitely presentable groups, and Turing machines. Those broad comparisons are motivational and are not used as universal theorems about the listed subjects.

The source also cites parity as evidence against a collapse to a simple theory: it says that certain ``irreducibly odd'' graph-links are stable in the sense that every equivalent graph contains the initial graph as a subgraph. The precise parity, irreducibility, and subgraph conventions are absent. This stability claim is retained as supplied context, not as a proof of undecidability or a substitute for the decision problem.
\subsection{Short English statement}
Decide whether equivalence of two finite graph-link diagrams admits an algorithm that always halts. The source conjectures undecidability and gives Turing-machine and parity-based motivation, but supplies no proof or exact graph-move presentation.
\subsection{Sources}
\begin{itemize}
\item Local UnsolvedMath registry, permanent ID 10600053, source code catalogue entry, virtual-knot problem 53. No individual proposer is supplied; the statement calls the negative answer ``our conjecture.''
\item \emph{Unsolved Problems in Virtual Knot Theory and Combinatorial Knot Theory}, revised survey, problem-list item 53. Supplied URL: \url{https://arxiv.org/abs/1409.2823}; recorded access 2026-07-29; extraction source-tex. Its dedicated section contains 60 top-level enumerated items with compound items retained intact.
\item The supplied summary lists ``Graph-Links'' (2010), \texttt{arXiv:1001.xxxx}; ``Introduction to Graph-Link Theory'' (2008), \texttt{arXiv:0810.xxxx}; ``Khovanov homology of graph-links'' (2010), \texttt{arXiv:1005.xxxx}; and ``Checkerboard graph links and simply laced Dynkin diagrams'' (2019), \texttt{arXiv:1907.xxxx}. The placeholder identifiers are incomplete and are not verified citations.
\item Supplied status partially\_solved and PARTIAL-PROGRESS; the summary describes graph-link theory as developed, says the decision question appears open, and reports no progress on undecidability found. These are editorial assessments. Background requires a current primary-literature status audit and redistribution review of the publicly accessible author/publisher copy, both NEEDS\_REVIEW. The registry assigns L4 whereas background records a default L3 assignment.
\end{itemize}
\end{document}
